\section{Histogram đơn giản [Trung bình]}

\subsection{Phân tích \& Thiết kế (M0)}
\begin{itemize}
    \item \textbf{Input:}
    \begin{itemize}
        \item File \texttt{array.txt} chứa: số nguyên $n$ và mảng $n$ số thực.
        \item File \texttt{config.txt} chứa: số nguyên $k$ là số lượng bin.
    \end{itemize}

    \item \textbf{Xử lý:}
    \begin{itemize}
        \item \textbf{Ý tưởng:}
        \begin{itemize}
            \item Chia đều khoảng $[\min, \max]$ thành $k$ bin.
            \item Độ rộng mỗi bin:
            \[
            w = \frac{\max - \min}{k}
            \]
            \item Với mỗi phần tử $x$, chỉ số bin:
            \[
            idx = \left\lfloor \frac{x - \min}{w} \right\rfloor
            \]
            \item Nếu $x = \max$ thì ép vào bin cuối cùng.
        \end{itemize}

        \item \textbf{Case biên:}
        \begin{itemize}
            \item $n = 0$ hoặc $k \le 0$ $\rightarrow$ file output rỗng.
            \item Nếu $\max = \min$ $\rightarrow$ đưa toàn bộ phần tử vào bin đầu tiên, các bin còn lại bằng 0.
        \end{itemize}

        \item \textbf{Module 1 (Đọc dữ liệu):} Hàm \texttt{readArray} và \texttt{readK}. Bắt lỗi mở file (\texttt{return false}).
        \item \textbf{Module 2 (Tìm min, max):} Hàm \texttt{findRange}.
        \item \textbf{Module 3 (Đếm bin):} Hàm \texttt{countBins}.
        \item \textbf{Module 4 (Ghi output):} Hàm \texttt{writeHistogram} ghi ra file \texttt{histogram.txt}.
    \end{itemize}

    \item \textbf{Output:}
    File \texttt{histogram.txt}:
\begin{lstlisting}
[lo, hi): count
\end{lstlisting}
\end{itemize}

\subsection{Mã nguồn C++11 (Tách module)}
\begin{lstlisting}
    #include <iostream>
    #include <fstream>
    #include <vector>
    #include <iomanip>
    #include <cmath>

    using namespace std;

    // MODULE 1a: Doc mang tu file
    bool readArray(string path, vector<double>& a) {
        ifstream file(path);
        if (!file.is_open()) {
            cerr << "Khong mo duoc file " << path << "\n";
            return false;
        }

        int n;
        file >> n;

        double x;
        for (int i = 0; i < n; i++) {
            file >> x;
            a.push_back(x);
        }

        file.close();
        return true;
    }

    // MODULE 1b: Doc so bin k
    bool readK(string path, int& k) {
        ifstream file(path);
        if (!file.is_open()) {
            cerr << "Khong mo duoc file " << path << "\n";
            return false;
        }

        file >> k;
        file.close();
        return true;
    }

    // MODULE 2: Tim min va max
    void findRange(const vector<double>& a, double& mn, double& mx) {
        mn = mx = a[0];
        for (double x : a) {
            if (x < mn) mn = x;
            if (x > mx) mx = x;
        }
    }

    // MODULE 3: Dem so phan tu trong moi bin
    vector<int> countBins(const vector<double>& a, int k, double mn, double mx) {
        vector<int> counts(k, 0);

        if (mx == mn) {
            counts[0] = (int)a.size();
            return counts;
        }

        double width = (mx - mn) / k;

        for (double x : a) {
            int idx = (int)((x - mn) / width);
            if (idx >= k) idx = k - 1;
            counts[idx]++;
        }

        return counts;
    }

    // MODULE 4: Ghi histogram.txt
    bool writeHistogram(string path, int k, double mn, double mx, const vector<int>& counts) {
        ofstream file(path);
        if (!file.is_open()) {
            cerr << "Khong ghi duoc file " << path << "\n";
            return false;
        }

        file << fixed << setprecision(2);

        double width = (mx - mn) / k;
        for (int i = 0; i < k; i++) {
            double lo = mn + i * width;
            double hi = (i == k - 1) ? mx : (mn + (i + 1) * width);
            file << "[" << lo << ", " << hi << "): " << counts[i] << "\n";
        }

        file.close();
        return true;
    }

    // MAIN
    int main() {
        vector<double> a;
        int k;

        if (!readArray("array.txt", a)) return 1;
        if (!readK("config.txt", k)) return 1;

        ofstream testCreate("histogram.txt");
        if (!testCreate.is_open()) {
            cerr << "Khong ghi duoc file histogram.txt\n";
            return 1;
        }
        testCreate.close();

        if (a.empty() || k <= 0) {
            return 0;
        }

        double mn, mx;
        findRange(a, mn, mx);

        vector<int> counts = countBins(a, k, mn, mx);

        if (!writeHistogram("histogram.txt", k, mn, mx, counts)) return 1;

        return 0;
    }
\end{lstlisting}

\subsection{Kế hoạch kiểm thử (Test Cases)}

\textbf{Test Case 1 (Thông thường)}

\textbf{Input (array.txt):}
\begin{lstlisting}
5 1.0 2.5 3.0 4.8 5.0
\end{lstlisting}

\textbf{Input (config.txt):}
\begin{lstlisting}
2
\end{lstlisting}

\textbf{Expected Output (histogram.txt):}
\begin{lstlisting}
[1.00, 3.00): 2
[3.00, 5.00): 3
\end{lstlisting}

\vspace{0.5cm}

\textbf{Test Case 2 (Case biên -- max = min)}

\textbf{Input (array.txt):}
\begin{lstlisting}
4 5 5 5 5
\end{lstlisting}

\textbf{Input (config.txt):}
\begin{lstlisting}
3
\end{lstlisting}

\textbf{Expected Output (histogram.txt):}
\begin{lstlisting}
[5.00, 5.00): 4
[5.00, 5.00): 0
[5.00, 5.00): 0
\end{lstlisting}

\vspace{0.5cm}

\textbf{Test Case 3 (Case biên -- n = 0)}

\textbf{Input (array.txt):}
\begin{lstlisting}
0
\end{lstlisting}

\textbf{Input (config.txt):}
\begin{lstlisting}
4
\end{lstlisting}

\textbf{Expected Output (histogram.txt):}
\begin{lstlisting}
\end{lstlisting}